\chapter{Assembler plusieurs photos}
L’onglet \emph{Assembler}, en haut du panneau de droite, crée une nouvelle image à
partir des photos cochées dans le ruban. Les photos d’origine ne changent pas. Le
calcul tourne en tâche de fond et peut être annulé ; vous pouvez continuer à parcourir
vos images pendant ce temps.

\section{Empiler pour réduire le bruit}
L’empilement moyenne plusieurs vues de la même scène, à la même exposition. Le bruit
baisse comme la racine du nombre de vues : quatre vues le divisent par deux, seize par
quatre. Les détails restent, car ils sont les mêmes d’une vue à l’autre.

\begin{enumerate}
\item Cochez au moins deux photos dans le ruban.
\item Choisissez la \emph{Référence} : l’image finale garde son cadrage. La photo du
milieu est proposée, car elle reste la plus proche des autres.
\item Cliquez sur \emph{Assembler}.
\item Regardez l’aperçu, choisissez le format, puis \emph{Exporter}.
\item \emph{Ouvrir le résultat} l’amène dans l’éditeur, pour le retoucher.
\end{enumerate}

Les vues sont recalées automatiquement : un trépied n’est pas nécessaire. Le résultat
est recadré sur la zone commune à toutes les vues. Ce qui bouge d’une vue à l’autre,
comme un passant, est écarté s’il n’apparaît que dans une minorité de vues. Une vue
impossible à recaler ou trop différente en exposition est écartée, et le message le
signale.

\section{Choisir le format}
\begin{itemize}
\item \emph{DNG} : le signal de l’appareil, déjà corrigé de l’objectif, sans profil
appliqué. Le fichier se rouvre comme un RAW : profil, balance et réglages restent libres.
\item \emph{TIFF} : l’image rendue en 16 bits avec le profil de la référence, lisible
partout.
\end{itemize}
Un fichier existant n’est jamais remplacé : le nom reçoit un numéro.

\section{Assembler un panorama}
Choisissez \emph{Panorama} en haut de l’onglet. Les photos prises du même endroit forment
une seule vue large, sur une rangée, horizontale ou verticale, ou sur plusieurs rangées.
Cochez-les dans le ruban, dans n’importe quel ordre, puis cliquez sur \emph{Assembler}.

Faites se recouvrir les photos voisines d’un quart à un tiers, et gardez les mêmes réglages
de prise. Une photo qui ne recouvre aucune autre, ou prise à une autre focale, est écartée, et
le message le signale.

\begin{itemize}
\item \emph{Le bas de la vue} sert d’appui : il reste de niveau, et le recadrage coupe plutôt
par le haut.
\item \emph{Projection} : \emph{Automatique} choisit selon le champ couvert. \emph{Plane}
garde les lignes droites sur un champ étroit, \emph{Cylindrique} convient à un long balayage,
\emph{Sphérique} à un champ large dans les deux sens.
\item \emph{Bords} : un panorama a rarement des bords droits. \emph{Recadrer} garde le plus
grand rectangle de pixels réels ; c’est le choix par défaut, rien n’est inventé. \emph{Remplir}
garde toute la surface et complète les coins avec des pixels inventés à partir du reste de
l’image. \emph{Laisser vides} garde les coins vides, à recadrer ensuite dans l’onglet Géométrie.
\end{itemize}
Le remplissage est convaincant sur le ciel. Sur la roche ou la végétation, regardez le résultat
à 100~\% : des blocs peuvent être mal placés. Le message de fin indique la part de pixels
inventés.

À main levée, un premier plan proche se déplace d’une photo à l’autre : les raccords passent
autant que possible là où cela ne se voit pas, mais tournez plutôt autour de l’appareil que
de vous-même.

\section{Photographier la Lune}
Un petit sujet sur un fond vide, comme la Lune, reçoit un traitement de plus. Les
déformations de l’atmosphère sont corrigées sur le sujet, et seule la moitié la plus
nette des vues y est fusionnée. Le reste du ciel garde toutes les vues. Le message
indique quand cette correction a eu lieu. Sur nos mesures à 200~mm, c’est surtout le
choix des vues nettes qui fait le gain : plus la série est longue, plus il y a de bonnes
vues parmi lesquelles choisir.

\section{Limites}
\begin{itemize}
\item Les réglages de prise doivent rester les mêmes : l’empilement ne fait pas de HDR.
\item Un objet proche, photographié à main levée, peut garder un bord un peu mou : il
ne se déplace pas comme le fond.
\item Le calcul demande environ 400~Mo de disque temporaire par photo de 32~Mpx pour un
empilement, et le double pour un panorama.
\item Un panorama ne déforme pas ses bords pour remplir le cadre.
\item Un profil de couleur pour le boîtier est nécessaire.
\end{itemize}
